% !TeX program = xelatex
\documentclass[a4paper,12pt]{article}

% ------------------------------------------------------------
% Idioma, fonte e formatação da página
% ------------------------------------------------------------
\usepackage[brazil]{babel}
\usepackage{fontspec}
\setmainfont{Times New Roman}

\usepackage[a4paper,margin=2.54cm]{geometry}
\usepackage{setspace}
\onehalfspacing
\usepackage{indentfirst}

% Títulos de seções e subseções em 12 pt e negrito
\usepackage{sectsty}
\allsectionsfont{\normalfont\normalsize\bfseries}

% Fonte exclusiva da capa.
% Para usar 12 pt também na capa, substitua a definição por:
% \newcommand{\fonteCapa}{\rmfamily\normalsize}
\newcommand{\fonteCapa}{\rmfamily\fontsize{14}{17}\selectfont}

% ------------------------------------------------------------
% Matemática, figuras e tabelas
% ------------------------------------------------------------
\usepackage{mathtools}
\usepackage{empheq}
\usepackage{systeme}
\usepackage{siunitx}
\usepackage[version=4]{mhchem}

\usepackage{graphicx}
\usepackage{xcolor}
\usepackage{tikz}
\usepackage{float}
\usepackage{subcaption}
\usepackage{booktabs}
\usepackage{adjustbox}
\usepackage{array}
\usepackage{tabularx}

\setlength{\marginparwidth}{2cm}
\usepackage[colorinlistoftodos]{todonotes}

% ------------------------------------------------------------
% Referências bibliográficas
% ------------------------------------------------------------
\usepackage[backend=biber,style=numeric,sorting=none]{biblatex}
\addbibresource{referencias.bib}
% ------------------------------------------------------------
% Listagens de código
% ------------------------------------------------------------
\usepackage{listings}

\definecolor{codegreen}{rgb}{0,0.6,0}
\definecolor{codegray}{rgb}{0.5,0.5,0.5}
\definecolor{codepurple}{rgb}{0.58,0,0.82}
\definecolor{backcolour}{rgb}{0.95,0.95,0.92}

\lstdefinestyle{mystyle}{
    backgroundcolor=\color{backcolour},
    basicstyle=\rmfamily\normalsize,
    commentstyle=\color{codegreen},
    keywordstyle=\color{magenta},
    numberstyle=\normalfont\normalsize\color{codegray},
    stringstyle=\color{codepurple},
    breakatwhitespace=false,
    breaklines=true,
    captionpos=b,
    keepspaces=true,
    numbers=left,
    numbersep=5pt,
    showspaces=false,
    showstringspaces=false,
    showtabs=false,
    tabsize=2,
    literate=
        {á}{{\'a}}1 {à}{{\`a}}1 {â}{{\^a}}1 {ã}{{\~a}}1
        {Á}{{\'A}}1 {À}{{\`A}}1 {Â}{{\^A}}1 {Ã}{{\~A}}1
        {é}{{\'e}}1 {è}{{\`e}}1 {ê}{{\^e}}1 {ë}{{\"e}}1
        {É}{{\'E}}1 {Ê}{{\^E}}1
        {í}{{\'i}}1 {î}{{\^i}}1 {Í}{{\'I}}1 {Î}{{\^I}}1
        {ó}{{\'o}}1 {ô}{{\^o}}1 {õ}{{\~o}}1
        {Ó}{{\'O}}1 {Ô}{{\^O}}1 {Õ}{{\~O}}1
        {ú}{{\'u}}1 {ù}{{\`u}}1 {û}{{\^u}}1 {Ú}{{\'U}}1
        {ç}{{\c{c}}}1 {Ç}{{\c{C}}}1 {°}{{\textdegree}}1
}
\lstset{style=mystyle}

% ------------------------------------------------------------
% Links e siglas
% ------------------------------------------------------------
\usepackage{hyperref}
\hypersetup{
    colorlinks=true,
    linkcolor=blue,
    filecolor=magenta,
    urlcolor=cyan,
    pdftitle={Projeto de pesquisa de doutorado},
    pdfauthor={Antonio Victor de Souza Correa}
}
\urlstyle{same}

\usepackage[acronym]{glossaries}
\makeglossaries

\newacronym{md}{MD}{\textit{Membrane Distillation}}
\newacronym{agmd}{AGMD}{\textit{Air Gap Membrane Distillation}}
\newacronym{vmd}{VMD}{\textit{Vacuum Membrane Distillation}}
\newacronym{dcmd}{DCMD}{\textit{Direct Contact Membrane Distillation}}
\newacronym{sgmd}{SGMD}{\textit{Sweep Gas Membrane Distillation}}
\newacronym{pp}{PP}{\textit{Polipropileno}}
\newacronym{pvdf}{PVDF}{\textit{Fluoreto de polivinilideno}}
\newacronym{ptfe}{PTFE}{\textit{Politetrafluoretileno}}

% ------------------------------------------------------------
% Dados da capa
% ------------------------------------------------------------
\newcommand{\nomeCandidato}{Antonio Victor de Souza Correa}
\newcommand{\areaConcentracao}{Termociências e Engenharia Térmica}
\newcommand{\tituloProjeto}{%
    Prevenção e recuperação de membranas poliméricas encrustadas
    de módulos de destilação por membrana com espaçamento de ar
    através de ultrassom e purga com ar comprimido%
}

\begin{document}

% ------------------------------------------------------------
% Capa
% ------------------------------------------------------------
\hypersetup{pageanchor=false}

\begin{titlepage}
    \thispagestyle{empty}
    \fonteCapa

    \begin{tikzpicture}[
        remember picture,
        overlay,
        every node/.style={inner sep=0pt,outer sep=0pt},
        texto/.style={text width=\textwidth,align=center}
    ]
        % Logos alinhadas às margens de 2,54 cm
        \node[anchor=north west]
            at ([xshift=2.54cm,yshift=-2.54cm]current page.north west)
            {\includegraphics[height=20mm]{ufrj-logo.pdf}};

        \node[anchor=north east]
            at ([xshift=-2.54cm,yshift=-2.54cm]current page.north east)
            {\includegraphics[height=20mm]{coppe-logo.pdf}};

        % A largura definida em "texto" permite a quebra automática.
        \node[texto,anchor=north]
            at ([yshift=-71mm]current page.north)
            {Projeto de pesquisa para o curso de doutorado no Programa
             de Pós-Graduação em Engenharia Mecânica da Universidade
             Federal do Rio de Janeiro};

        \node[texto]
            at ([yshift=-132mm]current page.north)
            {Candidato: \nomeCandidato};

        \node[texto]
            at ([yshift=-151mm]current page.north)
            {Área de concentração: \areaConcentracao};

        \node[texto]
            at ([yshift=-194mm]current page.north)
            {Título: \tituloProjeto};

        \node[texto,anchor=south]
            at ([yshift=2.54cm]current page.south)
            {Rio de Janeiro\\Dezembro de 2026};
    \end{tikzpicture}
    \null
\end{titlepage}

% ------------------------------------------------------------
% Sumário
% ------------------------------------------------------------
\pagenumbering{roman}
\hypersetup{pageanchor=true}
\tableofcontents
\thispagestyle{empty}

% ------------------------------------------------------------
% Corpo do projeto — Times New Roman 12 pt, espaçamento 1,5
% ------------------------------------------------------------
\clearpage
\pagenumbering{arabic}

\section{Introdução}

A água de produção é o principal rejeito da industria de óleo e gás, com geração diária de aproximadamente 39,75 milhões de metros cúbicos por dia~\cite{Ibrahim2023ProducedWater,Harkouss2024MembraneBased}, a qual é considerada em geral como salina ou hipersalina, com cloreto variando entre 36 a 210.224~$mgL^{-1}$~\cite{Alley2011ProducedWater}.

A destilação por membranas (MD) é um processo térmico movido pela diferença de pressão de vapor entre as superfícies da membrana, o qual é capaz de operar de forma estável com 1.000~$ppm$ de óleo e até 150.000~$ppm$ de sais, enquanto o processo de osmose reversa perde desempenho ao tratar soluções com 2.000~$ppm$ de sais e 150~$ppm$ de óleo~\cite{Mansi2023ProducedWater}.

O processo de destilação por membrana varia em função do modo como o permeado é recuperado. Nos módulos de destilação por membrana com espaçamento de ar, o vapor que permeia pelos poros da memebrana é condensado em uma placa de resfriamento e conduzido para um reservatório~\cite{Olatunji2018HeatMassTransport}.

Um dos obstáculos para a aplicação em campo é a queda de desempenho acarretada pela ocorrência de \textit{fouling} durante o tratamento de água de produção real, o que reduz o fluxo de permeado e auxilia a ocorrência de \textit{wetting}. Observa-se na literatura o uso de diversas técnicas para a restauração do fluxo, como pré tratamentos, aplicação de produtos químicos e limpeza da superfície com água, de modo a restaurar parcial ou completamente o fluxo de permeado. Entretanto, tais procedimentos consomem água, interrompem a operação ou precisam da compra, transporte e armazenamento de produtos químicos.

Estre projeto propõe avaliar 

\section{Relevância do tema}

\section{Objetivos do projeto}

\section{Metodologia}

\section{Viabilidade do projeto}

% Para exibir as referências citadas no projeto, descomente:
\printbibliography[title={Referências bibliográficas}]

\end{document}